% Part: proof-theory
% Chapter: proof-search
% Section: completeness

\documentclass[../../../include/open-logic-section]{subfiles}

\begin{document}

\olfileid[id]{pt}{ps}{cpl}

\olsection{Kelengkapan}

Sekarang kita menunjukkan bahwa algoritma pencarian bukti tidak mungkin
gagal: jika algoritma itu batal atau tidak pernah berhenti, tidak ada
!!{proof} bagi sekuen awal~$\Gamma \Sequent \Delta$. Untuk melakukannya,
kita mendefinisikan !!a{structure}~$\Struct M$ sedemikian sehingga bagi
semua $!A \in \Gamma$, $\Sat{M}{!A}$, dan bagi semua~$!A \in \Delta$,
$\Sat/{M}{!A}$. Dengan kata lain, semua !!{formula}s dalam~$\Gamma$ benar
dan semua !!{formula}s dalam~$\Delta$ salah dalam~$\Struct M$, sehingga
$\Gamma \Sequent \Delta$ tidak valid. Karena \Log{G3c} sahih,
$\Gamma \Sequent \Delta$ tidak mempunyai !!{proof}.

Kita mendefinisikan $\Struct{M}$ dari pasangan $\Theta, \Xi$ yang terdiri
atas !!{formula}s dan suatu himpunan~$C$ dari !!{constant}s. Kita memperoleh
$\Theta$, $\Xi$, dan~$C$ dari suatu \emph{cabang kegagalan} pada pelaksanaan
pencarian bukti yang gagal (batal atau tidak berhenti). Cabang kegagalan
adalah barisan sekuen $\Pi_n \Sequent \Lambda_n$ dan himpunan
!!{constant}s~$C_n$ sedemikian sehingga $\Pi_0 \Sequent \Lambda_0$ adalah
sekuen akhir $\Gamma \Sequent \Delta$, dan
$\Pi_{n+1} \Sequent \Lambda_{n+1}$ adalah suatu premis dari
$\Pi_n \Sequent \Lambda_n$ yang dihasilkan pada tahap~$n+1$ dan yang
untuknya algoritma tidak menemukan !!a{proof}. Premis seperti itu jelas harus
ada bagi setiap~$n$, sebab jika tidak, algoritma telah menemukan
!!a{proof}. Kita mengambil $C_n$ sebagai himpunan !!{constant}s yang
dikaitkan dengan~$\Pi_n \Sequent \Lambda_n$.

Jika algoritma batal pada sekuen paling atas yang hanya terdiri atas
!!{formula}s atomik, cabang yang berakhir pada sekuen itu merupakan cabang
kegagalan. Jika algoritma tidak pernah berhenti, algoritma terus menghasilkan
pohon yang semakin besar. Pohon-pohon ini dapat dipandang sebagai tahapan
pertumbuhan suatu pohon tak hingga (yakni pohon yang diperoleh setelah
menjalankan langkah algoritma sebanyak tak hingga). Pohon itu tak hingga,
tetapi bercabang berhingga (setiap simpul hanya mempunyai satu atau dua
anak---sekuen-sekuen yang tepat berada di atasnya). Berdasarkan Lemma
K\H{o}nig, setiap pohon tak hingga yang bercabang berhingga harus mempunyai
cabang tak hingga. Cabang tak hingga seperti itu merupakan cabang kegagalan
ketika algoritma pencarian berjalan selamanya.

Jika barisan $\Pi_n \Sequent \Lambda_n$ dan $C_n$ merupakan cabang
kegagalan, tetapkan $\Theta = \bigcup_n \Pi_n$ dan
$\Xi = \bigcup_n \Lambda_n$ (dengan menghapus indeks dan kemunculan ganda
!!{formula}s), serta $C = \bigcup_n C_n$.

Sekarang kita siap mendefinisikan~$\Struct{M}$.
\begin{enumerate}
\item Domain $\Domain{M}$ adalah himpunan suku yang dapat dibentuk dari~$C$
dan !!{function}s dalam~$\Gamma \Sequent \Delta$, jika ada, tetapi tanpa
!!{variable}s.
\item Jika $c \in \Domain{M}$, maka $\Assign{c}{M} = c$.
\item Jika $f$ adalah !!{function} $m$-tempat dalam
$\Gamma \Sequent \Delta$, dan $t_1$, \dots, $t_m \in \Domain{M}$, maka
$\Assign{f}{M}(t_1, \dots, t_m) = f(t_1, \dots, t_m)$.
\item Jika $R$ adalah !!{predicate} $m$-tempat dalam
$\Gamma \Sequent \Delta$, maka
$\Assign{R}{M} = \Setabs{\tuple{t_1, \dots, t_m} \in
\Domain{M}^m}{R(t_1, \dots, t_m) \in \Theta}$.
\end{enumerate}
$\Struct{M}$ merupakan \emph{model suku}: domainnya adalah himpunan suku,
dan !!{constant}s serta !!{function}s ditafsirkan sedemikian sehingga
$\Value{t}{M} = t$. !!^{formula}s atomik dalam~$\Theta$ digunakan untuk
mendefinisikan $\Assign{R}{M}$ sedemikian sehingga
$\Sat{M}{R(t_1, \dots, t_m)}$ jika dan hanya jika
$R(t_1, \dots, t_m) \in \Theta$.

\begin{lem}\ollabel{lem:truth}
Bagi semua $!A \in \Theta \cup \Xi$:
\begin{enumerate}
  \item Jika $!A \in \Theta$, maka $\Sat{M}{!A}$.
  \item Jika $!A \in \Xi$, maka $\Sat/{M}{!A}$.
\end{enumerate}
\end{lem}

\begin{proof}
Lakukan induksi pada $\depth{!A}$.

Pertama, andaikan $!A$ atomik, yakni $!A \ident \lfalse$ atau
$!A \ident R(t_1, \dots, t_m)$.

Karena $\Pi_n \Sequent \Lambda_n$ merupakan cabang kegagalan, tidak ada
$\Pi_n$ yang memuat $\lfalse$ (jika ada, $\Pi_n \Sequent \Lambda_n$
merupakan aksioma). Berdasarkan definisi~$\Struct{M}$,
$\Sat{M}{R(t_1, \dots, t_m)}$ jika dan hanya jika
$R(t_1, \dots, t_m) \in \Theta$. Secara bersama-sama, kedua hal ini
memberikan: jika $!A \in \Theta$, maka $\Sat{M}{!A}$.

Jika $!A$ atomik dan $!A \in \Pi_n$, maka $!A \in \Pi_{n+1}$; dan jika
$!A \in \Lambda_n$, maka $!A \in \Lambda_{n+1}$. (Hal ini mengikuti cara
algoritma pencarian bukti menghasilkan sekuen penerus.) Jadi, jika
$!A \in \Theta \cap \Xi$, maka bagi suatu~$n$,
$!A \in \Pi_n \cap \Lambda_n$. Ini berarti
$\Pi_n \Sequent \Lambda_n$ merupakan aksioma, sedangkan cabang kegagalan
tidak dapat memuat aksioma. Jadi, berdasarkan konstruksi,
$!A \notin \Theta \cap \Xi$ bagi $!A$ atomik; akibatnya, jika
$!A \in \Xi$, maka $!A \notin \Theta$. Berdasarkan definisi~$\Struct{M}$,
hal ini menyiratkan bahwa jika $!A \in \Xi$, maka $\Sat/{M}{!A}$.

Untuk langkah induksi, andaikan (1) dan (2) berlaku bagi semua !!{formula}s
dengan !!{depth} lebih rendah daripada~$!A$. Kita membuktikan (1) dan (2)
bagi $!A$ dengan membedakan beberapa kasus.

Pertama, perhatikan bahwa ketika $!A^i$ muncul dalam
$\Pi_n \Sequent \Lambda_n$, maka $!A^i$ juga muncul dalam
$\Pi_{n+1} \Sequent \Lambda_{n+1}$, kecuali jika $i$ merupakan indeks
terkecil dalam $\Pi_n \Sequent \Lambda_n$. Karena indeks terkecil meningkat
pada sekuen paling atas yang ditambahkan pada setiap tahap, pada akhirnya
kita mencapai $\Pi_n \Sequent \Lambda_n$ tempat $!A^i$ muncul dan $i$
merupakan indeks terkecil. Pada tahap $n+1$, algoritma mereduksi~$!A^i$.
Dengan kata lain, jika $!A \in \Theta$, maka bagi suatu $n$,
$\Pi_n = \Pi_n', !A^i$ dan $i$ merupakan indeks terkecil. Jika
$!A \in \Xi$, maka bagi suatu $n$,
$\Lambda_n = \Lambda_n', !A^i$ dan $i$ merupakan indeks terkecil.

\begin{enumerate}
  \item $!A \in \Theta$ dan $!A \ident !B \land !C$. Maka bagi suatu
  $n$, $\Pi_n = \Pi_n', (!B \land !C)^i$.
  $\Pi_{n+1} \Sequent \Lambda_{n+1}$ adalah premis bersesuaian dari
  \LeftR{\land}, yakni
  $\Pi_{n+1} = \Pi_n', !B^k, !C^{k+1}$. Akibatnya,
  $!B, !C \in \Theta$. Berdasarkan hipotesis induksi,
  $\Sat{M}{!B}$ dan $\Sat{M}{!C}$. Berdasarkan definisi $\Sat{M}{}$,
  kita memperoleh $\Sat{M}{!B \land !C}$.

  \item $!A \in \Xi$ dan $!A \ident !B \land !C$. Maka bagi suatu $n$,
  $\Lambda_n = \Lambda_n', !B \land !C$.
  $\Pi_{n+1} \Sequent \Lambda_{n+1}$ adalah salah satu dari dua premis
  \RightR{\land} dengan kesimpulan
  $\Pi_n \Sequent \Lambda'_n, !B \land !C$, yakni
  $\Lambda_{n+1} = \Lambda_n', !B^k$ atau
  $\Lambda_{n+1} = \Lambda_n', !C^k$. Akibatnya,
  $!B \in \Xi$ atau $!C \in \Xi$. Berdasarkan hipotesis induksi,
  $\Sat/{M}{!B}$ atau $\Sat/{M}{!C}$. Berdasarkan definisi $\Sat{M}{}$,
  kita memperoleh $\Sat/{M}{!B \land !C}$.

  \item Kasus $!A \ident \lnot !B$, $!A \ident !B \lor !C$, dan
  $!A \ident !B \lif !C$ serupa dan dibiarkan sebagai latihan.

  \item $!A \in \Theta$ dan $!A \ident \lexists[x][!B(x)]$. Maka bagi
  suatu $n$, $\Pi_n = \Pi_n', \lexists[x][!B(x)]^i$.
  $\Pi_{n+1} \Sequent \Lambda_{n+1}$ adalah premis bersesuaian dari
  \LeftR{\lexists}, yakni $\Pi_{n+1} = \Pi_n', !B(c)^k$ bagi suatu~$c$.
  Akibatnya, $!B(c) \in \Theta$ dan~$c \in C$. Berdasarkan hipotesis
  induksi, $\Sat{M}{!B(c)}$. Berdasarkan definisi $\Sat{M}{}$, kita
  memperoleh $\Sat{M}{\lexists[x][!B(x)]}$.

  \item $!A \in \Xi$ dan $!A \ident \lexists[x][!B(x)]$. Maka bagi suatu
  $n$, $\Lambda_n = \Lambda_n', \lexists[x][!B(x)]^i$. Setiap kali
  $\lexists[x][!B(x)]$ direduksi, $\lexists[x][!B(x)]$ tetap berada di sisi
  kanan sekuen premis dengan indeks baru. Karena itu,
  $\lexists[x][!B(x)]$ direduksi tak hingga kali di sisi kanan pada cabang
  kegagalan jika formula tersebut muncul di sana. Setiap suku
  $t \in \Domain{M}$ akhirnya menjadi ``suku pertama yang
  hanya memuat konstanta dalam~$C_n$ dan belum digunakan dalam suatu reduksi
  $\lexists[x][!B(x)]$ di sisi kanan'' pada cabang kegagalan. Jadi, bagi
  setiap $t \in \Domain{M}$, $!B(t) \in \Lambda_n$ bagi suatu~$n$, dan
  akibatnya $!B(t) \in \Xi$. Berdasarkan hipotesis induksi,
  $\Sat/{M}{!B(t)}$ bagi semua $t \in \Domain{M}$. Berdasarkan definisi
  $\Sat{M}{}$, kita memperoleh $\Sat/{M}{\lexists[x][!B(x)]}$.

  \item Kasus $!A \ident \lforall[x][!B(x)]$ serupa dan dibiarkan sebagai
  latihan.
\end{enumerate}
\end{proof}

\begin{cor}\ollabel{cor:complete} Algoritma pencarian bukti bagi
\Log{G3c} lengkap: jika algoritma itu batal atau tidak pernah berhenti,
sekuen awal~$\Gamma \Sequent \Delta$ tidak valid dan karenanya tidak
mempunyai !!{proof}.
\end{cor}

\begin{proof}
  Jika algoritma batal atau tidak pernah berhenti, terdapat cabang kegagalan
  yang darinya $\Struct{M}$ dapat didefinisikan. Berdasarkan
  \olref{lem:truth}, bagi semua $!A \in \Gamma$, $\Sat{M}{!A}$, dan bagi
  semua $!A \in \Delta$, $\Sat/{M}{!A}$. (Ingat bahwa
  $\Pi_0 = \Gamma$ dan $\Lambda_0 = \Delta$, sehingga
  $\Gamma \subseteq \Theta$ dan $\Delta \subseteq \Xi$.) Jadi,
  $\Gamma \Sequent \Delta$ tidak valid. Berdasarkan kesahihan~\Log{G3c},
  $\Gamma \Sequent \Delta$ tidak mempunyai !!{proof}.
\end{proof}

\begin{cor}
\Log{G3c} lengkap: jika $\Gamma \Sequent \Delta$ valid, maka
$\Log{G3c} \Proves \Gamma \Sequent \Delta$.
\end{cor}

\begin{proof}
  Kita membuktikan kontraposisinya. Andaikan $\Log{G3c} \Proves/
  \Gamma \Sequent \Delta$. Algoritma pencarian bukti harus batal atau tidak
  pernah berhenti karena tidak ada !!{proof} yang dapat ditemukan.
  Berdasarkan \olref{cor:complete}, $\Gamma \Sequent \Delta$ tidak valid.
\end{proof}

\end{document}
